\section{Theory and Method}

\begin{figure}[t]
  \centering
  \includegraphics[width=0.98\linewidth]{fig_method_overview.pdf}
  \caption{\textbf{EVQ-Cosh frequency allocation.} Left: geometric RoPE distributes frequencies uniformly in log-space. Center: the EVQ-cosh warp $\phi_k(\tau)$ bends the allocation below the identity (dashed), concentrating resolution at low frequencies; $\tau{\to}0$ recovers geometric RoPE. Right: the resulting allocation spreads low-frequency channels, retaining more positional information at long range.}
  \label{fig:method-overview}
\end{figure}

The derivation has three layers with distinct epistemic status:
\begin{enumerate}
\item \textbf{Exact (conditional on surrogate):} the ODE, closed-form family, geometric limit, and waterbed inequality (\S\ref{sec:ode}--\S\ref{sec:waterbed}).
\item \textbf{Practical branch choice:} drop the forcing term, retain the pure-tether warp (\S\ref{sec:evq-practical}).
\item \textbf{Semi-analytic calibration:} the $\tau^*$ scaling law, where one constant is empirically calibrated (\S\ref{sec:tau-star}).
\end{enumerate}
We separate these layers explicitly so that the exact variational results are not conflated with the practical recipe.

\subsection{RoPE as frequency allocation}

Let $\phi \in [0,1]$ parameterize log-frequency and define $\omega(\phi) = b^{-\phi}$. A discrete RoPE design is equivalent to choosing a density $\rho(\phi)$ over this interval and discretizing it by inverse CDF. Geometric RoPE corresponds to uniform $\rho$; EVQ derives $\rho$ from a variational objective. This reparameterization lifts the design space from a hand-written index rule to a continuous allocation problem.

\subsection{Exact kernel and broadband surrogate}

Under a distance prior $D(\Delta)$, the phase-collision kernel is
\begin{equation}
K(\phi_1,\phi_2)
=
\int D(\Delta)\cos(\omega(\phi_1)\Delta)\cos(\omega(\phi_2)\Delta)\,d\Delta.
\end{equation}
Our only approximation is a broadband surrogate,
\begin{equation}
K_{\mathrm{app}}(\phi_1,\phi_2)
\approx
\alpha \delta(\phi_1-\phi_2) + \beta \min(\phi_1,\phi_2),
\end{equation}
which isolates a diagonal ridge term and a smooth off-diagonal covariance. The $\delta$ term captures local ridge structure; $\min(\phi_1,\phi_2)$ is the Green kernel of a one-dimensional second-order operator. The exact kernel is oscillatory and cannot be pointwise-matched by a smooth two-parameter model; $K_{\mathrm{app}}$ captures the \emph{functional} structure of the collision objective. Appendix~\ref{sec:surrogate-validation} validates this: the EVQ solution derived from $K_{\mathrm{app}}$ reduces the collision score under the exact kernel by $24$--$92\%$ across 12 configurations spanning three orders of magnitude in $L$ and $b$. After this step, the remaining derivation is exact.

\subsection{Variational objective and exact ODE}
\label{sec:ode}

The broadband functional is
\begin{equation}
\mathcal{J}[\rho]
=
\frac{\alpha}{2}\int \rho(\phi)^2\,d\phi
+
\frac{\beta}{2}\int\!\!\int \rho(\phi)\rho(\psi)\min(\phi,\psi)\,d\phi\,d\psi
-
\mu\int \rho(\phi)b^{-2\phi}\,d\phi.
\end{equation}

\begin{theorem}[Stationary ODE]
Conditional on $K_{\mathrm{app}}$, the stationary density satisfies
\begin{equation}
\rho''(\phi) - \tau^2 \rho(\phi) = \gamma\, b^{-2\phi},
\qquad
\tau = \sqrt{\beta/\alpha},
\end{equation}
with a closed-form solution consisting of a hyperbolic tether and a Fisher-driven forcing term.
\end{theorem}

The ODE and its closed-form family are exact conditional on $K_{\mathrm{app}}$. The decomposition separates a global tether (redistributing density) from a local Fisher-utility forcing term, ensuring EVQ is the stationary compromise rather than an unconstrained long-range distortion.

\subsection{Closed-form EVQ and the geometric limit}

The pure-tether branch yields an inverse-CDF warp
\begin{equation}
\phi_k(\tau)
=
1 - \frac{1}{\tau}
\operatorname{arcsinh}\!\left((1-u_k)\sinh\tau\right),
\qquad
u_k = \frac{k}{K},
\end{equation}
and inverse frequencies $b^{-\phi_k(\tau)}$.

\begin{theorem}[Geometric limit]
EVQ-cosh converges smoothly to geometric RoPE as $\tau \to 0$. Geometric RoPE is therefore the $\tau{=}0$ degenerate limit of the EVQ family, not a generic optimum.
\end{theorem}

\subsection{Practical EVQ instantiation}
\label{sec:evq-practical}

For a head dimension $d_{\mathrm{head}}$ and training length $L$, the complete implementation is
\begin{equation}
\tau = \frac{d_{\mathrm{head}}}{\sqrt{L}},
\qquad
u_k = \frac{k}{K},
\qquad
\omega_k^{\mathrm{EVQ}} = b^{-\phi_k(\tau)},
\qquad
K=\frac{d_{\mathrm{head}}}{2}.
\end{equation}
The construction is deliberately minimal: set $\tau$, compute quantiles, evaluate the warp, replace inverse frequencies. No auxiliary loss, no extra parameters, no optimizer dependence. Two properties follow from the quantile form: $u_0{=}0$ implies $\phi_0(\tau){=}0$, anchoring the first channel at $\omega_0{=}1$; and the warp grows toward the low-frequency end, so EVQ preserves short-range structure without a separate design choice.

\begin{center}
\fbox{%
\parbox{0.93\linewidth}{%
\textbf{Practical recipe: EVQ-cosh initialization.}
Given $d_{\mathrm{head}}$, $L_{\mathrm{train}}$, and base $b$: (1) set $\tau=d_{\mathrm{head}}/\sqrt{L_{\mathrm{train}}}$; (2) define $K=d_{\mathrm{head}}/2$ and $u_k=k/K$; (3) compute
\[
\phi_k(\tau)=1-\frac{1}{\tau}\operatorname{arcsinh}\!\left((1-u_k)\sinh\tau\right);
\]
(4) initialize inverse frequencies as $b^{-\phi_k(\tau)}$. This is the complete method. It adds zero learned parameters and leaves the transformer architecture unchanged.}}
\end{center}

EVQ changes only the inverse-frequency initialization---distinct from learnable PE approaches that rely on in-distribution training dynamics to discover the allocation through extra parameters.

\subsection{Waterbed inequality}
\label{sec:waterbed}

\begin{proposition}
Under the broadband model, EVQ obeys a waterbed inequality: improving long-range frequency resolution necessarily redistributes error rather than eliminating it globally.
\end{proposition}

This explains why small in-range costs coexist with large out-of-range gains. Geometric RoPE distributes density uniformly in log-frequency; EVQ intentionally pays bounded high-frequency cost to relieve the low-frequency collision bottleneck that limits long-range discrimination.

\subsection{The \texorpdfstring{$\tau^*$}{tau*} scaling law: semi-analytic calibration}
\label{sec:tau-star}

\emph{Epistemic status: this subsection crosses from exact variational results into a semi-analytic layer. The purpose is practical: derive a $\tau^*$ rule that is good enough for deployment, not one that is exact.}

The surrogate provides $\sqrt{\beta/\alpha} \propto \sqrt{d_{\mathrm{head}}}$ but predicts only $L^{-0.11}$, too weak by a factor of~5. The missing $L$-dependence originates from the softmax attention map, and incorporating it requires one empirically calibrated constant $\lambda$.

\begin{proposition}[Softmax Transport Temperature Selection]
\label{prop:softmax-transport}
Define $\mathcal{F}(\tau) = S_{\chi^2}(\tau) - \lambda\,U(\tau, L)$,
where $S_{\chi^2}(\tau)$ is the Pearson~$\chi^2$ stiffness and $U(\tau,L)$ is the post-softmax transport utility (see Appendix~\ref{sec:stiffness-derivation} for definitions). The $1/L$ factor arises from the softmax Jacobian at diffuse baseline $p_0 = 1/L$.
Then the minimizer satisfies $\tau_* \propto d_{\mathrm{head}}/L^{\gamma}$; exact numerical optimization yields $\gamma = 0.465$ for Pearson~$\chi^2$ stiffness (Appendix~\ref{sec:stiffness-derivation}, Table~\ref{tab:stiffness-sweep}), which we round to $\gamma \approx 0.5$ for the practical rule.
\end{proposition}

The $\chi^2$ stiffness reflects attention-load asymmetry: diluted channels ($\rho < 1$) bear per-channel cost ${\propto}1/\rho$. KL, Jeffreys, and all R\'{e}nyi divergences produce $L$-exponents far from $-0.5$; only the $\chi^2$ family succeeds (Table~\ref{tab:stiffness-sweep}). The resulting scaling law
\begin{equation}
\tau^*(L) = \frac{d_{\mathrm{head}}}{\sqrt{L}}
\end{equation}
is semi-analytic: $d_{\mathrm{head}}$ dependence and $\gamma$ are derived, but $\lambda$ is calibrated. The formula drops the forcing branch and absorbs one free constant from data.

\paragraph{Sub-optimality as robustness.} We regard this deliberate sub-optimality not as a theoretical deficit but as the paper's central practical finding. A single surrogate-derived rule, operating without access to the exact kernel or per-setting tuning, matches the per-configuration empirical optimum to within $<1\%$ PPL across 99 runs spanning 27 configurations: from-scratch text (50M--750M), MLA with 16 channels, video DiT, and post-hoc LoRA at 7B (Appendix~\ref{sec:tau-scaling}). This implies a wide robustness basin: the frequency allocation landscape is locally flat near the variational solution, so even a rough analytic approximation lands inside the basin. The practical consequence is that a zero-parameter initialization outperforms learned positional encoding (Table~\ref{tab:pe-dominant}) despite being theoretically imperfect. Closing $\lambda$ analytically is an open question, but the observed basin width suggests the remaining headroom is small.

\subsection{Predictions}

The waterbed inequality predicts bounded short-range cost with large long-range gains, and a base dependence: EVQ helps most when the collision bottleneck is large (high $b$, short $L$); dead-zone experiments confirm this (Appendix). EVQ and inference-time scaling address orthogonal deficiencies: EVQ corrects \emph{shape} (non-uniform density), YaRN corrects \emph{range} (Figure~\ref{fig:orthogonal}). The corrections are additive, predicting the strong composition observed in \S\ref{sec:experiments}. Attention distance measurements on a trained 750M model verify the waterbed crossover at ${\approx}508$ tokens (Figure~\ref{fig:attn-mechanism} in Appendix).

\begin{figure}[t]
  \centering
  \includegraphics[width=0.90\linewidth]{fig_unification_orthogonal.pdf}
  \caption{\textbf{EVQ and YaRN correct orthogonal deficiencies of geometric RoPE.} (a)~EVQ bends the allocation curve \emph{below} the diagonal (shape correction), while YaRN shifts it \emph{above} (range correction). (b)~EVQ+YaRN combines both corrections additively. Setup: base${}=10{,}000$, $d_{\mathrm{head}}{=}128$, $L_{\mathrm{train}}{=}4096$, $4\times$ extrapolation.}
  \label{fig:orthogonal}
\end{figure}
